\documentclass[12pt]{article}
\usepackage[utf8]{inputenc}
\usepackage[T1]{fontenc}
\usepackage{hyperref}
\title{Tandaganon Language Resources for K-12 Curriculum Development}
\author{Compiled by Claude Sonnet 5 (claude-sonnet-5) for Andr\'as Kornai}
\date{}
\begin{document}
\maketitle
\section{Introduction}
Tandaganon (ISO 639-3 tgn, Glottocode tand1258) is an Austronesian language of central Surigao del Sur, Philippines, around Tandag, Tago, San Miguel, Bayabas and Cagwait. It has about 100,000 speakers. It is closely related to Surigaonon. Speakers call it Tandaganon or Tagon-on. Very little educational material was found.
\section{Primary Educational Resources}
No primers, textbooks or school materials in Tandaganon were found.
\section{Linguistic and Reference Materials}
\begin{itemize}
\item Glottolog entry for the Surigao family. It records the ISO split that gave Tandaganon its own code and the naming notes. \url{https://glottolog.org/resource/languoid/id/suri1274}
\item Wikipedia on Surigaonon. It discusses Tandaganon as a separate language or a southern variety. \url{https://en.wikipedia.org/wiki/Surigaonon_language}
\end{itemize}
\section{Religious and Community Resources}
\begin{itemize}
\item Joshua Project profile of the Tandaganon people. \url{https://joshuaproject.net/people_groups/20544}
\end{itemize}
\section{Digital and Audio Resources}
\begin{itemize}
\item A web page presenting a poem in Tandaganon with a translation. \url{https://phspirits.com/idlivun-tandaganon-translation}
\end{itemize}
\section{Curriculum Development Context}
No dictionary, grammar, orthography or school programme for Tandaganon was found. Curriculum work would start from the closely related Surigaonon and from field documentation.
\end{document}